\item \textbf{Datación de una roca lunar mediante Rubidio-Estroncio.}
En una muestra de roca traída de la Luna, se mide que contiene $1.82 \times 10^{10}\text{ átomos}$ de $^{87}\text{Rb}$ por gramo de material y $1.07 \times 10^9\text{ átomos}$ de $^{87}\text{Sr}$ por gramo. La desintegración relevante es $^{87}\text{Rb} \to ^{87}\text{Sr} + e^- + \bar{\nu}$, con una vida media de $4.75 \times 10^{10}\text{ años}$.
\begin{enumerate}
    \item Calcule con precisión la antigüedad geológica de la roca lunar.
    \item ¿Qué suposición implícita fundamental se realiza sobre las condiciones iniciales al aplicar este método de datación radiactiva?
\end{enumerate}